\documentclass[11pt]{scrartcl} % Font size

\usepackage{amsmath, amsfonts, amsthm, amssymb} % Math packages

\usepackage{pdfpages} %add pdf pages

\usepackage{float} %figures in place

\usepackage{physics,mathtools} %extra math

\usepackage{listings} % Code listings, with syntax highlighting

\usepackage{xcolor, soul}

\usepackage[english]{babel} % English language hyphenation

%\usepackage[xetex]{graphicx}
\usepackage{graphicx} % Required for inserting images
\graphicspath{{Figures/}{./}} % Specifies where to look for included images (trailing slash required)

\usepackage{booktabs} % Required for better horizontal rules in tables

\numberwithin{equation}{section} % Number equations within sections (i.e. 1.1, 1.2, 2.1, 2.2 instead of 1, 2, 3, 4)
\numberwithin{figure}{section} % Number figures within sections (i.e. 1.1, 1.2, 2.1, 2.2 instead of 1, 2, 3, 4)
\numberwithin{table}{section} % Number tables within sections (i.e. 1.1, 1.2, 2.1, 2.2 instead of 1, 2, 3, 4)

\setlength\parindent{0pt} % Removes all indentation from paragraphs

\usepackage{enumitem} % Required for list customisation
\setlist{noitemsep} % No spacing between list items


%----------------------------------------------------------------------------------------
%   HYPER LINKS
%----------------------------------------------------------------------------------------
\usepackage{hyperref} %for hyperlinks
\hypersetup{
    colorlinks=true,
    linkcolor=blue,
    filecolor=magenta,
    urlcolor=cyan,
}
\usepackage{extarrows} %xleftrightarrow

%----------------------------------------------------------------------------------------
%	DOCUMENT MARGINS
%----------------------------------------------------------------------------------------

\usepackage{geometry} % Required for adjusting page dimensions and margins

\geometry{
	paper=a4paper, % Paper size, change to letterpaper for US letter size
	top=2.5cm, % Top margin
	bottom=3cm, % Bottom margin
	left=3cm, % Left margin
	right=3cm, % Right margin
	headheight=0.75cm, % Header height
	footskip=1.5cm, % Space from the bottom margin to the baseline of the footer
	headsep=0.75cm, % Space from the top margin to the baseline of the header
	%showframe, % Uncomment to show how the type block is set on the page
}

%----------------------------------------------------------------------------------------
%	FONTS
%----------------------------------------------------------------------------------------

\usepackage[utf8]{inputenc} % Required for inputting international characters
\usepackage[T1]{fontenc} % Use 8-bit encoding

% \usepackage{fourier} % Use the Adobe Utopia font for the document

%----------------------------------------------------------------------------------------
%	SECTION TITLES
%----------------------------------------------------------------------------------------

\usepackage{sectsty} % Allows customising section commands
\sectionfont{\vspace{6pt}\centering\normalfont\scshape} % \section{} styling
\subsectionfont{\normalfont\bfseries\centering} % \subsection{} styling
\subsubsectionfont{\normalfont\itshape\centering} % \subsubsection{} styling
\paragraphfont{\normalfont\scshape} % \paragraph{} styling

%----------------------------------------------------------------------------------------
%	HEADERS AND FOOTERS
%----------------------------------------------------------------------------------------

\usepackage{scrlayer-scrpage} % Required for customising headers and footers

\ohead*{} % Right header
\ihead*{} % Left header
\chead*{} % Centre header

\ofoot*{} % Right footer
\ifoot*{} % Left footer
\cfoot*{\pagemark} % Centre footer
 % Include the file specifying the document structure and custom commands
% LaTeX settings for MATLAB code listings
% based on Ted Pavlic's settings in http://links.tedpavlic.com/ascii/homework_new_tex.ascii
\usepackage{listings}
% \usepackage[usenames,dvipsnames]{color}
\usepackage{color}

% This is the color used for MATLAB comments below
\definecolor{MyDarkGreen}{rgb}{0.0,0.4,0.0}

% For faster processing, load Matlab syntax for listings
\lstloadlanguages{Matlab}%
\lstset{language=Matlab,                        % Use MATLAB
        frame=single,                           % Single frame around code
        basicstyle=\scriptsize\ttfamily,        % Use small true type font
        keywordstyle=[1]\color{blue}\bfseries,  % MATLAB functions bold and blue
        keywordstyle=[2]\color{purple},         % MATLAB function arguments purple
        keywordstyle=[3]\color{blue}\underbar,  % User functions underlined and blue
        identifierstyle=,                       % Nothing special about identifiers
                                                % Comments small dark green courier
        commentstyle=\usefont{T1}{pcr}{m}{sl}\color{MyDarkGreen}\small,
        stringstyle=\color{purple},             % Strings are purple
        showstringspaces=false,                 % Don't put marks in string spaces
        tabsize=3,                              % 5 spaces per tab
        %
        %%% Put standard MATLAB functions not included in the default
        %%% language here
        morekeywords={xlim,ylim,var,alpha,factorial,poissrnd,normpdf,normcdf,imresize},
        %
        %%% Put MATLAB function parameters here
        morekeywords=[2]{on, off, interp},
        %
        %%% Put user defined functions here
        morekeywords=[3]{FindESS, homework_example, my_imfilter},
        %
        morecomment=[l][\color{blue}]{...},     % Line continuation (...) like blue comment
        numbers=left,                           % Line numbers on left
        firstnumber=1,                          % Line numbers start with line 1
        numberstyle=\tiny\color{blue},          % Line numbers are blue
        stepnumber=1                            % Line numbers go in steps of 5
        }

% Includes a MATLAB script.
% The first parameter is the label, which also is the name of the script
%   without the .m.
% The second parameter is the optional caption.
\newcommand{\matlabscript}[2]
  {\begin{itemize}\item[]\lstinputlisting[caption=#2,label=#1]{#1.m}\end{itemize}}


\usepackage{fontspec}
\setmainfont{Tinos Nerd Font} %nice font for english and greek

\begin{document}

\begin{titlepage}
    \centering
    \includegraphics[width=0.9\textwidth]{logo_tuc.png}\par\vspace{1cm}
    \normalfont\normalsize
    \textsc{\textcolor[rgb]{0.66, 0.09, 0.19}{Ηλεκτρολόγων Μηχανικών και Μηχανικών Υπολογιστών}}\\ % Your university, school and/or department name(s)
    \vspace{25pt} % Whitespace
    %\textcolor[rgb]{0.66, 0.09, 0.19}{\rule{\linewidth}{0.5pt}}
    \rule{\linewidth}{0.5pt}\\ % Thin top horizontal rule
    \vspace{20pt} % Whitespace
    {\Huge Μηχανική Όραση}\\ % The assignment title

    {\huge Αναφορά Τρίτου Project}\\ % The assignment title
    \vspace{12pt} % Whitespace
    \rule{\linewidth}{2pt}\\ % Thick bottom horizontal rule
    \vspace{12pt} % Whitespace
    \vspace{2cm}

    {\LARGE{Τσιαούσης Χρήστος \hfill Μητσάκης Χαράλαμπος}
        \par
        \texttt{2016030017 \hfill 2006030075}
        \par
    }

    \vfill
    Διδάσκουσα

    Ε. Δούτση

    \vfill

% Bottom of the page
    {\large \today\par}
\end{titlepage}

\newpage



\section{Εισαγωγή}

Η συγκεκριμένη εργασία διερευνά την αναγνώριση αντικειμένων σε μια εικόνα
βάσει της μεθοδολογίας Histogram of Oriented Gradients και Sliding Window Classification.

Αποτελείται από τα παρακάτω μέρη:
\begin{itemize}
  \item \textbf{Gradient Generation} Η συνάρτηση \texttt{mygradient} παράγει τις κλίσεις των
  magnitude και orientation για κάθε pixel της εικόνας εισόδου.
  \item \textbf{Histogram of Oriented Gradients} Η συνάρτηση \texttt{hog} παράγει τα orientation
  histograms για κάθε 8x8 superpixel της εικόνας εισόδου.
  \item \textbf{Template Detection} Η συνάρτηση \texttt{detect} δέχεται μία εικόνα εισόδου καθώς
  και κάποια template θετικών κι αρνητικών match. Επιστρέφει της συντεταγμένες και τα σκορ για
  τα καλύτερα matches που βρήκε.
\end{itemize}

\section{Αποτελέσματα}
Παρατηρήθηκε πως το sliding window classification λειτουργεί καλά αν τα feature που έχουν οριστεί
στο θετικό template του training set βρίσκονται σε παρόμοιο μέγεθος στο testing set. Δηλαδή, αν
και το HOG για τα block 8x8 παράγει orientation independent identities για κάθε block, πρέπει το
feature να κάνει καλό fit στις περιοχές που διαλέγονται στο train έτσι ώστε να γίνεται έλεγχος
ομοιότητας στο test έχοντας παρόμοια 8 block σε κάθε window (8 block των 8x8 δίνουν παράθυρο 64x64).

\subsection{Image Gradient}

Καλώντας την mygradient ως εξής:

\begin{verbatim}
mygradient(im2double(rgb2gray(imread('../data/test3.jpg'))));
\end{verbatim}

παίρνουμε τα παρακάτω αποτελέσματα:

\begin{figure}[H]
  \includegraphics[width=16cm]{../output/1_mag.jpg}
  \caption{magnitude}
\end{figure}

\begin{figure}[H]
  \includegraphics[width=16cm]{../output/1_ori.jpg}
  \caption{orientation}
\end{figure}

\begin{figure}[H]
  \includegraphics[width=16cm]{../output/1_hogDraw.jpg}
  \caption{HOG draw}
\end{figure}

\subsection{Detection}
Για το detection υπάρχει μία training εικόνα στην οποία διαλέγουμε 4 windows (μεγέθους
64x64) τα οποία θα είναι το θετικό template και βάσει αυτών θα ``ψάχνει'' matches
ο classifier.

\begin{figure}[H]
  \includegraphics[width=16cm]{../output/3_positive.jpg}
  \caption{positive template}
\end{figure}

Όμοια, διαλέγουμε και 4 παράθυρα για το αρνητικό template.

\begin{figure}[H]
  \includegraphics[width=16cm]{../output/3_negative.jpg}
  \caption{negative template}
\end{figure}

Παρακάτω φαίνεται πως η αναγνώριση είναι ανεξάρτητη orientation αλλά εξαρτημένη του scaling factor.

\begin{figure}[H]
  \includegraphics[width=16cm,trim={0 2.5cm 0 2.5cm},clip]{../output/3_result_1.jpg}
  \caption{faces\_1b.jpg}
\end{figure}

\begin{figure}[H]
  \includegraphics[width=16cm]{../output/3_result_2.jpg}
  \caption{faces\_1b.jpg rotated $90^{\circ}$}
\end{figure}

\begin{figure}[H]
  \includegraphics[width=16cm,trim={0 2.5cm 0 2.5cm},clip]{../output/3_result_3.jpg}
  \caption{faces\_1b.jpg scaled x2}
\end{figure}

\begin{figure}[H]
  \includegraphics[width=16cm]{../output/3_result_4.jpg}
  \caption{faces2.jpg}
\end{figure}

\section{Συμπεράσματα}

Φαίνεται η μέθοδος να μην επιρρεάζεται από το orientation των αντικειμένων αλλά
δεν αναγνωρίζει απόλυτα templates διαφορετικού μεγέθους. Για να αντιμετοπιστεί
το πρόβλημα αυτό θα μπορούσαν να εφαρμοστούν ποικίλες υλοποιήσεις. Για παράδειγμα,
κάποιος θα μπορούσε να δημιουργήσει ένα \textbf{array από scaling factors} έτσι
ώστε να δημιουργηθούν διάφορα training sets, όπου τα επιλεγμένα feature του χρήστη
επιλέγονται και σε διαφορετικά scales με αποτέλεσμα να γίνει η αναγνώριση invariant
από το μέγεθος του εκάστοτε feature στο testing set. Μια άλλη ιδέα θα ήταν να χρησιμοποιήσουμε
διαφορετική μέθοδο classification. Αφού τα labels είναι μόνο δύο (match/όχι match)
θα μπορούσαμε να υλοποιήσουμε εκπαίδευση με κάποιο SupportVectorMachine για καλύτερο
διαχωρισμό των κλάσεων.



\end{document}
